\documentclass[11pt]{article}
\usepackage[margin=2.6cm]{geometry}
\usepackage{parskip}
\usepackage[hidelinks]{hyperref}
\usepackage{booktabs}

% Title = the record's citable identity (author decision 2026-07-16):
% the tool, not the process. Matches CITATION.cff.

\title{\vspace{-1.2em}CycloneNet: A Reproducible Pipeline and Leakage-Safe
Two-Basin Dataset for Tropical-Cyclone Rapid-Intensification Analysis\\[0.3em]
\large Project Status --- Zenodo record 10.5281/zenodo.18751255, version 3.0.0}
\author{Estefano Senhor Ferreira \\
        Senior Software Engineer, Independent Researcher \\
        \texttt{estefano.senhor@gmail.com} \\
        \url{https://github.com/estefano-ferreira/cyclone-net}
}
\date{2026-07-16}

\begin{document}
\maketitle
\vspace{-1.5em}

\noindent This document states what the project is today (\S1--\S2) and
records, in full, the claims of version~2.0.0 that were refuted or
superseded (\S3). The version~2.0.0 manuscript follows below,
\emph{unmodified}, as the historical record of what was claimed --- read
it through \S3. Full audit trail: \texttt{ERRATA.md},
\texttt{BENCHMARK.md}, and \texttt{docs/hypothesis\_registry.md} in the
project repository
(\url{https://github.com/estefano-ferreira/cyclone-net}).

\section*{1. What this project is}

A two-basin (East Pacific + North Atlantic), 1980--2023, hurricane-season,
leakage-safe tropical-cyclone rapid-intensification (RI) dataset, with an
auditable construction pipeline:

\begin{itemize}
  \item \textbf{Population:} 16{,}780 events from 992 storms --- 799 RI
  positives, 15{,}962 negatives, and 19 events whose label is undefined
  (NULL) under the corrected strict-temporal labeling (dataset~v2; the
  v1$\to$v2 per-row provenance ships with the dataset).
  \item \textbf{Reproducible labeling:} the labeling chain is
  byte-reproducible from the raw IBTrACS source (32{,}989/32{,}989 rows,
  every label identical on reconstruction).
  \item \textbf{Leakage-safe splits:} storm-level, hash-deterministic,
  with a frozen benchmark override map; no storm crosses splits.
  \item \textbf{Lifecycle provenance:} raw ERA5 was processed in windows
  (download $\to$ extract $\to$ verify $\to$ discard), each with a
  SHA-256-checksummed provenance manifest; repairs to stored artifacts
  (basin metadata, label correction) carry their own manifests.
  \item \textbf{Pre-registration:} the modeling claims were subjected to a
  pre-registered campaign with four verdicts, each confidence interval
  read once (\S2).
  \item \textbf{No reference model:} the retired CNN's frozen test
  metrics are historical record; the tabular baseline was evaluated on
  development folds only and never on the frozen test set.
  \item \textbf{Licensing:} dual --- code under MIT, dataset under
  CC~BY~4.0 (mandatory Copernicus/ERA5 and IBTrACS attributions).
\end{itemize}

\section*{2. The pre-registered campaign (2026-07-16)}

All hypotheses were registered with decision rules before any result was
read; each confidence interval was read once, and the reports are
committed in the repository. The hypothesis registry
(\texttt{docs/hypothesis\_registry.md}) is what makes the verdicts below
traceable rather than asserted: every hypothesis --- tested, refuted,
cancelled or deferred --- appears there with its dated registration, the
decision rule fixed before the result, and the verdict; refuted entries
are preserved, not removed.

\begin{itemize}
  \item \textbf{H6 --- do shear and mid-level RH channels add skill?}
  NULL. Cross-seed mean $\Delta$PR-AUC $=+0.0185$, 95\%~CI
  $[-0.0070,\,+0.0431]$ includes zero. No detectable contribution at this
  resolution and regime, for this architecture.
  \item \textbf{H9/V1 --- CNN vs.\ tabular baseline (validity).} NEGATIVE.
  $\Delta_1=$ PR-AUC(CNN) $-$ PR-AUC(GBM$_{SF}$) $=-0.0781$, 95\%~CI
  $[-0.1162,\,-0.0422]$: a gradient-boosted model over 44 aggregate
  scalars beats the CNN.
  \item \textbf{H9/V2 --- does the architecture justify itself at a fixed
  information diet?} NULL. $\Delta_2=$ PR-AUC(CNN) $-$ PR-AUC(GBM$_F$)
  $=+0.0005$, 95\%~CI $[-0.0285,\,+0.0316]$. By pre-registered
  consequence, \textbf{the CNN architecture is retired}. Scope guard,
  verbatim from the registry: ``V2 speaks about THIS model, not about
  whether spatial structure carries information.''
  \item \textbf{H8 --- FuelMap physics-loss ablation.} CANCELLED:
  undecidable once V2 retired the architecture those losses shape. The
  pre-registration and harness remain in the repository as record.
\end{itemize}

Two boundaries these verdicts do \emph{not} cross: (i)~the H9 baseline is
a gradient-boosted model built from this project's own released data ---
comparison against operational baselines (e.g.\ SHIPS-RII) has \emph{not}
been performed and remains a precondition for any operational skill
claim; (ii)~per the scope guard, none of this establishes whether spatial
structure in the fields carries RI information --- only that this CNN did
not extract it.

\section*{3. Correction record: what version 2.0.0 claimed, and what testing found}

\textbf{Energy-source localization (FuelMap): REFUTED, not pending.}
Version~2.0.0 stated that external spatial validation against TCHP/OHC
``remains future work.'' It was executed and returned negative, from
three independent angles: (i)~\emph{static TCHP validation} ($n=226$
eligible test events): the FuelMap peak beats a random-point null
($p=0.0003$) but does \emph{not} locate the TCHP peak better than a naive
storm-centre baseline --- median 539~km vs.\ 561~km, closer in only 46\%
of events ($p=0.30$, sign-flip permutation); (ii)~a \emph{dynamic
displacement test} against a control; (iii)~a \emph{physics-prior
control}: the pure enthalpy-flux prior, with no learned weights,
reproduces the dynamic behavior --- it is arithmetic of the prior, not
learned skill. A counterfactual occlusion test shows the RI prediction
depends on the identified region, which demonstrates
\emph{model-internal} dependence only.

\textbf{Headline metrics: SUPERSEDED --- do not cite.} The reported
ROC-AUC~0.83 / recall~0.905 on 2{,}193 samples (211 RI positives) were
never reproducible from the public artifacts. The reproducible numbers,
from the released 1980--2023 dataset, are: ROC-AUC~0.796
[95\%~CI 0.753--0.837], PR-AUC~0.251 [0.179--0.331], recall~0.852, on a
frozen test split of 2{,}679 events (115 RI positives, v1 labels; see
\S1 for the v2 label correction). These are the metrics of a model that
has since been retired (\S2); they stand as historical record, not as a
reference model.

\textbf{``Physics-guided'': label RETIRED (2026-07-16).} The only
equation-consistency term in the released code is disabled
($\lambda_{\mathrm{consistency}}=0$) and documented as near-degenerate;
the remaining terms are weak supervision toward a heuristic prior whose
semantics were refuted (hypothesis H1). There are no conservation laws or
imposed dynamics. Class and module names in the code retain the
historical identifiers because the released checkpoint references them.

\textbf{``Atmospheric Singularity Mapping'': REMOVED as an overstatement}
(\texttt{ERRATA.md}, item~5). A literature check found the components and
the forensic framing established, not novel; the contribution of this
project is engineering and reproducibility, not a discovery.

\textbf{The CNN: RETIRED} by the pre-registered verdict H9/V2 (\S2).

The version~2.0.0 manuscript remains attached below without
modification, as the record of what was claimed.

\section*{4. Provenance of this document}

This project began as a software engineer's attempt to apply
critical-systems engineering discipline to a geophysical hypothesis. The
hypothesis was refuted by the external validation, and the model was
retired by verdicts the author pre-registered before reading any result.
What remains is the dataset, the audit trail, and the negative results
documented above.

\end{document}
